\documentclass[10pt]{article}
\usepackage[margin=21mm]{geometry}
\usepackage{amsmath,amssymb,amsthm,booktabs,hyperref}
\newtheorem{theorem}{Theorem}\newtheorem{lemma}[theorem]{Lemma}\newtheorem{proposition}[theorem]{Proposition}
\newcommand{\ii}{\mathrm i_2}\newcommand{\nuu}{\nu_2}\newcommand{\F}{\mathbb F}\newcommand{\GL}{\operatorname{GL}}\newcommand{\GU}{\operatorname{GU}}\newcommand{\Sp}{\operatorname{Sp}}\newcommand{\PSp}{\operatorname{PSp}}\newcommand{\PGSp}{\operatorname{PGSp}}\newcommand{\Aut}{\operatorname{Aut}}
\setlength{\emergencystretch}{2em}
\title{Explicit symplectic reverser and overgroup bounds}
\author{Human Compression project}\date{}
\begin{document}\maketitle
This mathematical component is for $G=\PSp_{2m}(q)$, $m\ge2$.
The Singer maximal-overgroup classification for $m\ge5$ is retained from
Bubboloni--Lucido--Weigel Lemma3.1, citing Malle--Saxl--Weigel Theorem1.1.
It includes a unitary overgroup omitted from the target's Table4. The
low-rank covering interface $m=2,3,4$ is explicitly pending verification;
the estimates below also cover the corresponding corrected candidate list.
All source finite exceptions are acquired from the finite35 certificates.
The lower bounds are acquired from the accepted classical lower theorem,
and the nonalternating isometry group calculation from the companion
linear/unitary proof. This proves the symplectic positive branch for $m\ge5$ relative to that
explicit classification premise; the low-rank cover and the complete original
disposition are not claimed here.
Write $\nuu(H)=\ii(H)+1$, the number of solutions of $h^2=1$.

\section{The torus and all possible reversers}
Put $E=\F_{q^{2m}}$, $E_0=\F_{q^m}$, and $\bar u=u^{q^m}$. Choose
$\gamma\ne0$ with $\bar\gamma=-\gamma$, and take the alternating form
\[
 B(u,v)=\operatorname{Tr}_{E/\F_q}(\gamma u\bar v).
\]
Multiplication by a generator $\omega$ of the norm-one group of $E/E_0$
is an isometry $X$ of order $q^m+1$. Its projective image $x$ has order
$M=(q^m+1)/\gcd(2,q-1)$. In $J=\PGSp_{2m}(q)$ its torus centralizer
has order $D=q^m+1$.

Here is the regularity and outer-coset check. For $m\ge2$ in the unexcluded
parameters, $M>q^{m/2}+1$. A projective centralizer normalizes the field
algebra generated by $X$ and induces a $q^j$-power map on $E$. It follows
that $M\mid q^j-1$. For $0<j<2m$, the gcd with $q^m+1$ is either at most
two or is $q^{\gcd(m,j)}+1$, where the gcd is a proper divisor of $m$.
It is at most $q^{m/2}+1$, a contradiction. Thus the centralizer consists
of multiplication maps, and its projective similitude part is the torus of
order $D$. The same argument applied after squaring shows that a natural
projective inverter induces exactly the $q^m$-power map. It therefore has a
lift $h:u\mapsto c\bar u$. This is a similitude precisely when
$a=c\bar c\in\F_q^*$; its multiplier is $-a$ and its square is $aI$.

For $m\ge3$, the quotient $\Aut(G)/J$ is the cyclic field group. A
nontrivial involutory field projection has $f$ even, $q=p^f$, and exponent
$f/2$. Eigenvalue comparison would imply
$M\mid p^{f/2+fj}+1$ for some $j$, and hence $M\mid q^{2j+1}-1$.
But $\gcd(q^m+1,q^{2j+1}-1)\le2$, impossible. The same excludes the
ordinary field involution for $m=2$. Thus all involutory reversers are
natural, except possibly the exceptional graph-field coset of $\Sp_4(q)$
in characteristic two. There are at most $D$ reversers in any one coset,
since their pairwise quotients centralize $x$.

For $\Sp_4(q)$ in characteristic two, the outer group is cyclic of order
$2f$. If $f$ is even its involution is a field automorphism and is excluded
above. If $f$ is odd, the additional involutory coset has fixed group
$\operatorname{Sz}(q)$, by Aschbacher--Seitz Section19, (19.3)--(19.5).
These exact outer-automorphism statements, with the 1978 erratum retained,
are borrowed structural inputs. We may uniformly count at most $2D$
reversers in degree four and characteristic two; all other cases use $D$.

\section{Their fixed groups and involutions}
Suppose first that $q$ is odd. Since $h^2=aI$ and its multiplier is $-a$,
$h$ is skew-adjoint for $B$. If $a$ is a square, its two $m$-dimensional
eigenspaces are paired totally isotropic spaces, and the isometric
centralizer is $\GL_m(q)$. If $a$ is nonsquare, the algebra $\F_q[h]$ is
$\F_{q^2}$ and the same pairing gives a hermitian form; the centralizer is
$\GU_m(q)$. Projective commutation permits a lift to send $h$ to $-h$,
so $C_G(\alpha)$ is contained in the corresponding normalizer, modulo
$\langle-I\rangle$, with quotient of order at most two over the linear or
unitary part. For $m\ge3$, its image on the derived projective group is a
subgroup of $\Aut(\operatorname{PSL}_m(q))$ or
$\Aut(\operatorname{PSU}_m(q))$. The kernel is cyclic of order at most
$q+1$. The extra coset acts as inverse transpose, a genuine outer graph
map in these degrees; it introduces no further kernel.
Lawther--Liebeck--Seitz Proposition1.3 consequently gives
\begin{equation}\label{oddmu}
 \ii(C_G(\alpha))\le (q+1)\{2(q^{a_m}+q^{a_m-1})+1\}-1,
 \qquad a_m=(m^2+m-2)/2.
\end{equation}
The entire kernel size is charged: it is not assumed central in the
normalizing coset.

For $q$ even, rescale a natural inverter to have square one. Hilbert90 for
$E/E_0$ gives the required conjugating multiplier. Multiplying that
multiplier by an element of $E_0$ makes its norm one, since squaring is
onto in $E_0$; thus the conjugation is by a torus isometry and takes the
inverter to $u\mapsto\bar u$.
Its square-zero part has rank $m$. The induced form is nonalternating,
since
$B(u,\bar u)=\operatorname{Tr}_{E/\F_q}(\gamma u^2)$ is not identically
zero. In a symplectic basis it is therefore
\[
 t=\begin{pmatrix}I&I\\0&I\end{pmatrix}.
\]
Every commuting isometry has diagonal blocks $A,A$, where $AA^t=I$, and
its off-diagonal block can be written $DA^{-t}$ for a symmetric matrix $D$.
Equivalently the centralizer is
\[
 \operatorname{Sym}_m(q)\rtimes J_m(q),
 \qquad A:D\longmapsto ADA^t,
\]
where $J_m$ is the isometry group of a nonalternating symmetric form.
This follows directly from the two block commutation equations and the
symplectic pairing equation; the kernel is the additive space of symmetric
matrices, of order $q^{m(m+1)/2}$. In particular
\begin{equation}\label{evenmu}
 \ii(C_G(\alpha))\le q^{m(m+1)/2}\nuu(J_m(q))-1.
\end{equation}
The identity image is included. The acquired elementary description gives
$J_m=\Sp_{m-1}$ for odd $m$ and
\[
 \nuu(J_m)\le q^{m-1}+q^{m-2}\ii(\Sp_{m-2})\quad(m\text{ even}).
\]
For $m=2$, $J_2$ is the additive group of order $q$. Each nonidentity
$A\in J_2$ fixes a two-dimensional subspace of the three-dimensional
space $\operatorname{Sym}_2(q)$: conjugate its single Jordan block to
$\left(\begin{smallmatrix}1&1\\0&1\end{smallmatrix}\right)$ and solve
$ADA^t=D$. Thus the exact centralizer involution count is
\begin{equation}\label{mu4even}
 \mu_2(q)=q^3+(q-1)q^2-1=2q^3-q^2-1.
\end{equation}
In the exceptional outer coset, the Suzuki group has one involution class,
with centralizer $q^2$, so its involution count is
$(q-1)(q^2+1)<\mu_2(q)$. Its class order is retained from the primary
involution inventory. This proves the same bound for every reverser.

\paragraph{The odd-characteristic degree-four case.}
Here the graph on the derived degree-two group is inner, so use an exact
count rather than the kernel assertion in~\eqref{oddmu}. In the split
case take
\[
 \GL_2(q): A\mapsto\operatorname{diag}(A,A^{-t}),\qquad
 w=\begin{pmatrix}0&I\\-I&0\end{pmatrix}.
\]
Then $w^2=-I$ and $w$ acts on $A$ by $A^{-t}$. In the nonsplit case the
same normalizer is described by $\GU_2(q)$ and a semilinear $w$ with square
$-I$, acting by conjugation on the field entries. Such a $w$ is obtained
by multiplying coordinate conjugation by a scalar of norm $-1$.
Modulo $\langle-I\rangle$, an element of the linear part is an involution
when $A^2=I$ or $-I$. Splitting their two eigenspaces, or using the
irreducible quadratic centralizer when needed, gives
\[
\begin{array}{c|cc}
& q\equiv1\pmod4&q\equiv3\pmod4\\\hline
\ii(\GL_2/\{\pm I\})&q(q+1)+1&q^2\\
\ii(\GU_2/\{\pm I\})&q^2&q(q-1)+1.
\end{array}
\]
In the other coset, square $\pm I$ means $A=\pm A^t$.
There are $q^2(q-1)$ nonsingular symmetric and $q-1$ nonsingular alternating
linear matrices, hence $(q-1)(q^2+1)/2$ projective solutions. For unitary
matrices, symmetric ones correspond by finite-field descent to the two
orthogonal forms, so their number is
$|\GU_2|(1/|O_2^+|+1/|O_2^-|)=q^2(q+1)$; there are $q+1$ alternating
ones. This gives $(q+1)(q^2+1)/2$ solutions. Consequently, for $q\ge5$,
\begin{equation}\label{mu4odd}
 \ii(C_G(\alpha))\le\tfrac12(q+1)(q^2+1)+q^2+q+1<q^3.
\end{equation}
The inequality reduces to $q^3-3q^2-3q-3>0$, already positive at five
and increasing thereafter.

\section{Retained covering and the uniform inequalities}
For $m\ge5$, retain the following complete list of maximal overgroup
\emph{types} from Bubboloni--Lucido--Weigel Lemma3.1 for the full Singer
lift $X$ in $\Sp_{2m}(q)$: the extension-field types
$\Sp_{2m/r}(q^r).r$ for primes $r\mid m$; additionally $O^-_{2m}(q)$
when $q$ is even; and additionally $\GU_m(q).2$ when $m,q$ are odd.
Their unitary notation uses the field $q^2$. The last type is absent from
the target Table4, so that printed list is not used as a complete cover.
Passing to projective groups causes no new case: the full inverse image
of a subgroup containing $x$ contains $X$ and the scalar centre.
The analogous list for $m=2,3,4$ is used only as an explicit conditional
interface until its small-rank classification has been verified.

Here are the multiplicities for the fixed element, rather than just the
conjugacy types. In an extension-field normalizer, $X^r$ is linear over
$K=\F_{q^r}$. Its order is $(q^m+1)/\gcd(r,q^m+1)$ and is greater than
$q^{m/2}+1$ in all unexcluded parameters. For $m\ge6$ use
$q^{m/2}-1>m$; for $m=2,3,4$ direct substitution at the smallest retained
$q$ proves the inequality. For $m=5$, the only remaining smallest case is
$q=2$, where $r=5$ does not divide $33$. The same eigenvalue gcd argument
as above proves that $X^r$ is irreducible. Thus
$K\subseteq C_{\operatorname{End}_{\F_q}(E)}(X^r)=E$, where a subfield of
this order is unique. The full normalizer of this one field structure is
one subgroup, so there is at most one extension-field term for each $r$.
For the unitary type, apply the same argument to $X^2$. Its unique
quadratic field is the subfield $\F_{q^2}\subset E$; since $m$ is odd,
the bar map restricts to its nontrivial involution. The alternating trace
form determines the hermitian form, and the full normalizer of this field
structure is the one unitary normalizer. In characteristic two, two
$X$-invariant quadratic forms polarizing to $B$ differ by the square of a
linear functional fixed by $X$. Irreducibility excludes a nonzero such
functional. Hence the orthogonal overgroup is unique as well.
The normalizer of $\langle X\rangle$ preserves each of the unique
prime-degree extension fields, so it is included even when $m$ is prime.
No uncharged normalizer term is discarded.

For the extension-field term, the actual scalar quotient embeds in the
automorphism group of $\PSp_{2m/r}(q^r)$. For odd $r$, its involutions lie
in the linear part, but the same upper bound suffices. The bound is
\[
 E_r<2\{q^{m^2/r+m}+q^{m^2/r+m-r}\}.
\]
The additional even-characteristic orthogonal term is bounded by
$2(q^{m^2}+q^{m^2-1})$; its isometry group embeds in
$\Aut(\Omega^-_{2m}(q))$. This includes $m=3$ through $D_3=A_3$ and
$m=2$ through $O^-_4(q)=\operatorname{SL}_2(q^2).2$.
For the additional odd-characteristic unitary normalizer, the projective
image acts on $\operatorname{PSU}_m(q)$ with cyclic kernel of order at most
$q+1$; its extra coset is a graph automorphism for $m\ge3$. Thus its entire
involution contribution is at most the right side of~\eqref{oddmu}.
All smaller groups used here are simple. Thus, with at most $m$ prime
terms, their total is at most $3m q^{m^2/2+m}$; for $m=3$ it is at most
$3q^6$. In characteristic two add $3q^{m^2}$ for the orthogonal term.

For odd $q$ and $m\ge3$, the reversing contribution from~\eqref{oddmu}
together with that unitary term when present is less than
$5q^{(m^2+3m)/2}$. Indeed, charge $(q^m+2)\mu$ uniformly, and use
$q^m+2\le(29/27)q^m$, $q+1\le4q/3$, and
$2(q^{a_m}+q^{a_m-1})+1\le3q^{a_m}$.
Relative to the acquired lower scale
$q^{m^2+m}$, at $m=3$ the total is at most
$5/q^3+3/q^6<1/4$ for $q\ge3$. For $m\ge4$ it is at most
\[
 5\,3^{-m(m-1)/2}+3m\,3^{-m^2/2}<1/4.
\]
Both summands decrease from $m=4$, where their sum is less than $0.009$.
The accepted quarter lower bound proves the strict counting inequality.

For $q$ even the accepted explicit class lower bound is stronger:
$\ii(G)>(2/3)q^{m^2+m}$. For odd $m\ge5$,~\eqref{evenmu} and LLS give
reversing contribution less than $5q^{(3m^2+6m-1)/4}$. Its ratio, with the
other terms, is at most
\[
 5\,2^{-(m-1)^2/4}+3\,2^{-m}+3m\,2^{-m^2/2}<2/3.
\]
At $m=5$ the sum is less than $0.41$ and all terms decrease. For even
$m\ge6$, the bound $\nuu(J_m)<4q^{m^2/4+m/2-2}$ follows from its displayed
formula and LLS. The ratio is instead at most
\[
 5\,2^{-(m^2/4-m+2)}+3\,2^{-m}+3m\,2^{-m^2/2}<2/3;
\]
at $m=6$ it is less than $0.204$, again decreasing. The numerical LLS
bound for the nonsimple $\Sp_4(2)$ is paid by its exact $75$ involutions;
no simplicity premise is used there.

For $m=3$ in characteristic two, $\nuu(J_3)=q^2$, so the reversing term
is $(q^3+1)(q^8-1)$. At $q\ge4$ the total ratio is bounded by
$(1+q^{-3})/q+2(1+q^{-1})/q^3+3/q^6<2/3$.
For $m=4$, the exact $\nuu(J_4)=q^4+q^3-q^2$ bounds the reversing term
by $5q^{18}$. At $q\ge4$ the total ratio is at most
$5/q^2+3/q^4+12/q^8<2/3$. The unused $q=2$ cases are finite35 inputs.

For $m=2$, odd $q\ge5$, the accepted class lower bound exceeds
$q^5(q-1)/2$. Equations~\eqref{mu4odd} and the single extension-field
bound give total less than
$q^5+2q^4+q^3+2q^2$, which is smaller already at $q=5$; after division by
$q^5$, the left side $(q-1)/2$ increases and the right side decreases.
For even $q\ge8$, use the exact class lower bound
$(q^2-1)(q^4-1)$, the at most $2D$ reversers, and two overgroup bounds.
Their difference is at least
\[
 q^6-4q^5-3q^4-4q^3-q^2+3>0.
\]
Divide by $q^6$ and drop its positive last term; the resulting lower bound
is positive at eight and increasing. The other degree-four parameters are
in finite35, with the nonsimple $\Sp_4(2)$ excluded from the original
positive simple-group claim.

An involution outside the covered proper subgroups generates $G$ with $x$.
Any reversing automorphism fixing this partner squares to the identity,
so all such obstructions have been included. The strict inequalities
therefore prove the symplectic existence assertion for $m\ge5$ relative to
the explicit retained covering classification; for $m=2,3,4$ this remains
conditional on the stated low-rank cover, without the source's unexpanded arithmetic
or its uniform single-inverter-class assertion.

\paragraph{Attribution and limits.}
The target is Leemans--Liebeck, arXiv1603.07713v3, Section4. The missing
unitary overgroup was identified by root from BLW Lemma3.1 before this
component was frozen. Earlier drafts and their arithmetic are retained. The projective
coset organization was suggested by root; the present symplectic trace-form,
field exclusion, exact degree-four fibres and inequalities are this author's
new derivations. AS1976 Section19 and its 1978 erratum supply the exceptional
outer fixed groups; BT Proposition2.3/Table5 supplies the Suzuki class count;
LLS2002 Proposition1.3 is the numerical upper theorem. The finite35 and
classical lower components retain their separate authorship and distinct
reviews. This is neither a proof of maximal subgroup classification nor a
whole-source admission.
\begin{thebibliography}{9}
\bibitem{LL} D. Leemans and M. W. Liebeck, \emph{Chiral polyhedra and finite
simple groups}, arXiv:1603.07713v3, Section4, Tables3--5.
\bibitem{LLS} R. Lawther, M. W. Liebeck and G. M. Seitz,
\emph{Fixed point ratios in actions of finite exceptional groups of Lie type},
Pacific J. Math. 205 (2002), Proposition1.3.
\bibitem{Ber} A. Bereczky, \emph{Maximal overgroups of Singer elements in
classical groups}, J. Algebra 234 (2000), 187--206. Classification retained
through the exact statement in [1], not an acquired full-text proof.
\bibitem{BLW} D. Bubboloni, M. S. Lucido and Th. Weigel,
\emph{2-Coverings of classical groups}, arXiv:1102.0660, Lemma3.1;
there citing G. Malle, J. Saxl and T. Weigel, \emph{Generation of classical
groups}, Geom. Dedicata 49 (1994), 85--116, Theorem1.1.
\bibitem{AS} M. Aschbacher and G. M. Seitz, \emph{Involutions in Chevalley
groups over fields of even order}, Nagoya Math. J. 63 (1976), 1--91,
Section19; with the authors' 1978 corrections.
\bibitem{BT} T. C. Burness and A. R. Thomas, \emph{On the involution fixity
of exceptional groups of Lie type}, arXiv:1707.03562v2, Proposition2.3/Table5.
\end{thebibliography}
\end{document}

